\documentclass[11pt]{article}
\usepackage[a4paper,margin=26mm]{geometry}
\usepackage{amsmath,amssymb,amsthm,array,longtable,microtype}
\usepackage[T1]{fontenc}
\usepackage[hidelinks]{hyperref}
\newtheorem{theorem}{Theorem}
\newtheorem{lemma}[theorem]{Lemma}
\newtheorem{corollary}[theorem]{Corollary}
\newcommand{\N}{\mathbb N}
\newcommand{\DS}[1]{D(#1)}
\newcommand{\cl}{\operatorname{cl}}
\newcommand{\hatset}[1]{\widehat{#1}}
\newcommand{\op}{\mathrm{op}}
\setlength{\emergencystretch}{2em}
\title{A ZFC extension of Brandsma--van Hartskamp--van Mill's example\par
for Feldman--Wilce's bounded-clique question}
\author{Haobo Ma\thanks{ChronoAI Pte Ltd, Singapore.}\and
Wenlin Zhang\thanks{National University of Singapore; ChronoAI Pte Ltd, Singapore.}}
\date{18 November 2026\\[-1pt]{\small CC BY 4.0}}
\begin{document}
\maketitle
\begin{abstract}
Extending the conditional example of Brandsma, van Hartskamp and van Mill
(1997, Example~2.3, pp.~28--29), we answer yes in ZFC to Feldman and
Wilce's Section~5, Question~2, p.~11 of arXiv:2608.13320v1: the
bounded-clique property is strictly weaker than the finite open chromatic
property, as witnessed by a deleted-square construction.
The space $D(\beta\N)=(\beta\N)^2\setminus\Delta$ has countably compact
countable power, so every symmetric irreflexive relation closed in its
whole square has a finite bound on clique size. Its coordinate-swap
relation nevertheless has no countable cover by open independent sets.
Clopen separation proves the compactness assertion, and an infinite
Ramsey argument proves the coloring obstruction. We also give an explicit
higher power that is not countably compact and derive the sphere lower
bound suggested by Feldman and Wilce from Holmsen--Lee's theorem.
\end{abstract}

\section{The question and its witness}
For a topological space $X$, a \emph{closed orthogonality} is a symmetric,
irreflexive relation $R\subseteq X^2$ closed in the entire product,
including its diagonal. An $R$-clique is a subset whose distinct points
are pairwise related. Following Feldman--Wilce~\cite{FW}, $X$ has the
\emph{bounded-clique property} if, for each closed orthogonality $R$,
there is an integer $N_R$ bounding the cardinality of every $R$-clique.
The bound depends on $R$; in particular every clique is finite.
The \emph{open chromatic number} $\chi_\op(R)$ is the least cardinality
of a cover of $X$ by open $R$-independent sets. Cover members may overlap.
The \emph{finite open chromatic property} requires
$\chi_\op(R)<\infty$ for every closed orthogonality $R$.

Compactness implies the finite open chromatic property, which implies
the bounded-clique property. Feldman--Wilce's Question~2 asks whether
the second implication is strict~\cite[Section~5, p.~11]{FW}.
We use the Stone space $K=\beta\N$ of ultrafilters on the discrete
natural numbers and its deleted square
\[
 X=\DS K=\{(p,q)\in K^2:p\ne q\},
 \qquad \tau(p,q)=(q,p).
\]
Both $X$ and its products carry their indicated subspace and product
topologies, respectively.

\begin{theorem}\label{main}
In ZFC the Hausdorff space $X=\DS{\beta\N}$ has countably compact
countable power and the bounded-clique property. The relation
\[
 R_\tau=\{(x,\tau(x)):x\in X\}\subseteq X^2
\]
is a closed orthogonality with clique number exactly $2$ and
$\aleph_1\le\chi_\op(R_\tau)\le 2^{\aleph_0}$. Consequently $X$ does not have the
finite open chromatic property.
\end{theorem}

An earlier conditional separation follows from Brandsma--van
Hartskamp--van Mill~\cite[Example~2.3, pp.~28--29]{BHM}.
They credit the involution construction in that example to LaBerge.
Under $\diamondsuit$ they describe a first-countable, countably compact,
perfectly normal space $E$ with a fixed-point-free involution and with
$\beta E$ its one-point compactification. First countability makes $E$
sequentially compact: an infinite sequence either repeats a point
infinitely often or has a cluster point, from whose countable local
base a convergent subsequence can be chosen. Hence Feldman--Wilce's
Corollary~4(a) gives the bounded-clique property. The involution extends
to $\beta E$ and fixes its unique remainder point. It therefore has no
finite coloring in the sense of~\cite[pp.~25--26]{BHM}. In that definition
colors are functionally open; here every open set is functionally open
because $E$ is perfectly normal. Its closed involution graph thus has
no finite open independent cover. Theorem~\ref{main} removes the
additional set-theoretic hypothesis by a different, explicit space.

The point-pair separation argument below is self-contained for compact
Hausdorff extremally disconnected spaces. A compact $F$-space is a compact
Hausdorff space in which disjoint cozero sets have disjoint closures;
a cozero set is $f^{-1}(\mathbb R\setminus\{0\})$ for a continuous
real-valued function $f$. Protasov--Pym--Strauss
\cite[Lemma~2.1, p.~211]{PPS} prove that unequal pairs in a compact
Hausdorff $F$-space have a common subsequence whose endpoint sets have
disjoint closures. In an extremally disconnected compact space these
compact sets have a clopen separator, giving Lemma~\ref{pair} below.
Their Two-Function Lemma~2.2 (p.~212), applied to the projections
$K^2\to K$, already gives countable compactness of $D(K)$ for every
compact Hausdorff $F$-space $K$. Thus the separation mechanism is prior
work; we retain a self-contained proof under the stronger
extremal-disconnectedness hypothesis.

\section{Clopen separation and countable compactness}
A compact Hausdorff space is \emph{extremally disconnected} if the
closure of every open set is open. Such a space has a clopen base:
given $x\in U$ open, regularity gives $x\in V\subseteq\cl V\subseteq U$,
and $\cl V$ is clopen. Moreover, disjoint open sets have disjoint
closures. Indeed, if $V,W$ are disjoint open sets, then
$W\cap\cl V=\varnothing$; the clopen set $\cl V$ also misses $\cl W$.

For $A\subseteq\N$, put
$\hatset A=\{p\in\beta\N:A\in p\}$. These sets form a clopen base,
and the principal ultrafilters identify $\N$ with a dense subspace.
If $O\subseteq\beta\N$ is open, then $O\cap\N$ is dense in $O$ and
$\cl(O\cap\N)=\hatset{O\cap\N}$. Thus $\beta\N$ is extremally
disconnected. We use the usual compact Hausdorff topology of this
Stone space.

\begin{lemma}\label{pair}
Let $K$ be compact Hausdorff and extremally disconnected. For any
sequence $(p_n,q_n)$ with $p_n\ne q_n$, there is a clopen
$C\subseteq K$ such that $p_n\in C$ and $q_n\notin C$ for infinitely
many indices $n$.
\end{lemma}
\begin{proof}
Suppose otherwise. Applying the supposition to $C$ and $K\setminus C$
shows that for every clopen $C$, the equivalence
$p_n\in C\Longleftrightarrow q_n\in C$ holds eventually.

We construct decreasing clopen sets $C_0=K,C_1,\ldots$, decreasing
infinite sets of available indices $L_n$, and strictly increasing
$k_n\in L_n$. Both endpoints of every indexed pair in $L_n$ lie in
$C_n$. Given $k_n$, choose a clopen $H\subseteq C_n$ containing
$p_{k_n}$ and missing $q_{k_n}$. After discarding finitely many
indices, the two endpoints of each remaining pair agree on membership
in $H$. Infinitely many pairs therefore lie either both in $H$ or
both in $C_n\setminus H$. Let $C_{n+1}$ be that side, and let
$L_{n+1}\subseteq L_n$ be the corresponding infinite set of indices
larger than $k_n$. Label the endpoints of the pair at $k_n$ as
$x_n,y_n$ so that
\[
 x_n\in D_n:=C_n\setminus C_{n+1},\qquad y_n\in C_{n+1}.
\]
The shells $D_n$ are pairwise disjoint and clopen. If $y_i\in D_j$,
then $i<j$.

For each $i$ let $r(i)$ be the unique shell index containing $y_i$,
if such an index exists, and put $r(i)=i$ otherwise. Define
$b_0=0$ and $b_{n+1}=1+\max(b_n,r(b_n))$. Then
\[
 y_{b_n}\notin D_{b_m}\quad\hbox{for all }n,m.
\]
For $m\le n$ this follows from the preceding shell order; for $m>n$
we have $b_m>r(b_n)$. Write $B=\{b_n:n\in\N\}$ and
$Q(A)=\cl\bigl(\bigcup_{i\in A}D_i\bigr)$ for $A\subseteq B$.
Splitting any infinite $A\subseteq B$ into two disjoint infinite parts
gives disjoint clopen sets $Q(A')$ and $Q(A'')$. At least one misses
any prescribed point.

Consequently we can choose decreasing infinite sets
$A_0=B,A_1,\ldots$ such that $y_{b_n}\notin Q(A_{n+1})$.
Choose strictly increasing $s_n\in A_n$ and put $S=\{s_n:n\in\N\}$.
For every $n$, the set $S\setminus A_{n+1}$ is finite. The identity
of the union with its two parts gives
\[
 Q(S)\subseteq Q(A_{n+1})\ \cup\!
 \bigcup_{i\in S\setminus A_{n+1}}D_i.
\]
The point $y_{b_n}$ misses both sets on the right. Thus the clopen
set $Q(S)$ contains $x_i$ and misses $y_i$ for every $i\in S$.
It distinguishes the original pair at $k_i$ infinitely often,
contrary to eventual agreement. This proves the lemma.
\end{proof}

\begin{lemma}\label{power}
If $K$ is compact Hausdorff and extremally disconnected, then
$\DS K^{\,\N}$, meaning $(\DS K)^\N$, is countably compact.
\end{lemma}
\begin{proof}
If $\DS K$ is empty the assertion is immediate. Otherwise let
$z_n=(z_n(j))_{j\in\N}$ be any sequence in $(\DS K)^\N$.
Apply Lemma~\ref{pair} to its zeroth coordinate. On an infinite
set of indices, that coordinate lies in $C_0\times(K\setminus C_0)$
for a clopen $C_0$. Repeat on this infinite set for coordinate $1$,
then coordinate $2$, and so on, producing nested infinite index sets.
A diagonal subsequence $z_{a_n}$ eventually satisfies each coordinate's
clopen constraint. Compactness of $(K^2)^\N$ supplies a cluster point
$w=(w(j))_j$ of this subsequence. Each closed rectangle
$C_j\times(K\setminus C_j)$ contains its $j$th coordinate eventually,
so contains $w(j)$. Thus $w(j)$ has unequal coordinates for every $j$,
and $w\in(\DS K)^\N$. The subspace topology makes $w$ a cluster
point there as well.

For completeness, in a $T_1$ space the sequence cluster-point property
implies countable compactness. If an increasing countable open cover
$V_n$ had no finite subcover, choose $u_n\notin V_n$. A cluster point
in $V_N$ would have a neighborhood containing no terms with $n\ge N$,
a contradiction. Conversely, countable compactness gives a cluster
point of every sequence by intersecting the decreasing nonempty
closed sets $\cl\{u_n:n\ge N\}$. This also covers sequences with
repetitions. Applying the first implication proves the assertion.
\end{proof}

The same countable-power conclusion holds for every compact Hausdorff
$F$-space $K$, as a consequence of~\cite[Lemma~2.1]{PPS}.
For each coordinate successively, thin to an infinite set of indices
whose two endpoint sets have disjoint closures $A_j,B_j\subseteq K$.
A diagonal subsequence eventually belongs to $A_j\times B_j$ in
coordinate $j$. Any cluster point in the compact space $(K^2)^\N$
then belongs to $(D(K))^\N$, since $A_j\cap B_j=\varnothing$ for every
$j$. This is a deduction from the cited separation lemma, rather than
a countable-power statement quoted from that paper.

\begin{lemma}[Feldman--Wilce, Theorem 1]\label{criterion}
If $Y^\N$ is countably compact, then $Y$ has the bounded-clique
property.
\end{lemma}
\begin{proof}
We give the argument of~\cite[Section~3, pp.~7--8]{FW}.
The empty space is immediate; otherwise projection shows that $Y$
is countably compact. An infinite clique would contain a sequence
of distinct points with cluster point $y$. For any neighborhoods
$U,V$ of $y$, infinitely many sequence indices lie in $U\cap V$;
two distinct terms therefore give an edge in $U\times V$.
Closedness of the relation would imply $(y,y)\in R$, contrary to
irreflexivity. All cliques are thus finite.

If finite clique sizes were unbounded, choose an $(n+1)$-clique
$\{y_{n,0},\ldots,y_{n,n}\}$ for every $n\ge0$, and define
$v_n(j)=y_{n,\min(j,n)}$. Let $v\in Y^\N$ be a cluster point of
the sequence $v_n$. For $i\ne j$ and neighborhoods $U$ of $v(i)$
and $V$ of $v(j)$, some $n\ge\max(i,j)$ satisfies
$v_n(i)\in U$, $v_n(j)\in V$. These are distinct members of one
clique, so $(v(i),v(j))\in R$ by closedness. Irreflexivity makes
the coordinates $v(i)$ pairwise distinct. They form an infinite
clique, contradicting the preceding paragraph.
\end{proof}

\section{The swap graph has no finite open coloring}
The graph $R_\tau$ is closed in $X^2$: it is the intersection of
the two closed equality sets asserting $p=s$ and $q=r$ for
$x=(p,q)$ and $y=(r,s)$. It is symmetric because $\tau^2$ is the
identity, and irreflexive because $p\ne q$. Each point has exactly
one neighbor. A clique therefore has at most two points, and
$\{x,\tau(x)\}$ attains that bound.

\begin{lemma}\label{color}
No finite family of open $R_\tau$-independent subsets covers $X$.
\end{lemma}
\begin{proof}
Suppose $U_1,\ldots,U_m$ is such a cover; $m=0$ already fails
because $X$ is nonempty. For distinct principal points $a,b\in\N$,
write
\[
 P(a,b)=\{i:(a,b)\in U_i\}\subseteq\{1,\ldots,m\}.
\]
Color the unordered pair $\{a,b\}$, where $a<b$, by the ordered
pair $(P(a,b),P(b,a))$. The infinite Ramsey theorem for finite
colorings of pairs gives an infinite $A\subseteq\N$ and fixed
patterns $P,Q$ such that
\[
 P(a,b)=P,\qquad P(b,a)=Q\quad(a<b,\ a,b\in A).
\]
Here $P\cap Q=\varnothing$, since otherwise a single $U_i$ would
contain a point and its swap.

The finite-color Ramsey fact used here has a short proof. Choose
$a_0$ and an infinite set of larger vertices whose pairs with $a_0$
have one fixed color. Choose $a_1$ in that set and thin its remaining
tail to one color from $a_1$, and continue. The infinite pigeonhole
principle gives an infinite subsequence of the $a_j$ whose chosen
colors agree; all pairs in that subsequence have that color.

Split $A$ into disjoint infinite sets $A_0,A_1$. Extend the cofinite
filter on each support to free ultrafilters $p,q$ containing $A_0,A_1$,
respectively. Their disjoint supports imply $p\ne q$. Choose
$i$ with $(p,q)\in U_i$. Openness and the clopen rectangle base give
$B_0\in p$, $C_0\in q$ such that every unequal pair in
$\hatset{B_0}\times\hatset{C_0}$ belongs to $U_i$. Put
$B=B_0\cap A_0$ and $C=C_0\cap A_1$. These are disjoint infinite
subsets of $A$, since members of free ultrafilters are infinite.
Every principal pair in $B\times C$ therefore belongs to $U_i$.
There exist $b\in B,c\in C$ with $b<c$, forcing $i\in P$.
There also exist $b'\in B,c'\in C$ with $c'<b'$, forcing $i\in Q$.
The two principal pairs need not be swaps; their common color
membership and the homogeneous patterns give the contradiction
$i\in P\cap Q$.
\end{proof}

Apply Lemmas~\ref{power} and~\ref{criterion} with $K=\beta\N$,
and then Lemma~\ref{color}. Together with the graph observations,
these prove the finite-color assertion of Theorem~\ref{main}.
Since $X$ is countably compact, a countable open cover would have a finite
subcover; hence $\chi_\op(R_\tau)\ge\aleph_1$.
For the upper bound, the $2^{\aleph_0}$ open sets
$\hatset A\times\hatset{\N\setminus A}$, $A\subseteq\N$, cover $X$
and are independent. Distinct ultrafilters are separated by one such $A$.
This completes Theorem~\ref{main}. All choices above are within ZFC.

\section{Consequences and source coverage}
\begin{corollary}\label{products}
For every nonempty finite or countably infinite set $I$, the
Hausdorff space $X^I$ has the bounded-clique property and fails
the finite open chromatic property.
\end{corollary}
\begin{proof}
The space $(X^I)^\N$ is homeomorphic to $X^{I\times\N}$, hence
to $X^\N$, so Lemma~\ref{criterion} applies. Fix $i_0\in I$.
Pulling $R_\tau$ back along the coordinate projection defines a
closed orthogonality on $X^I$. If it admitted a finite open independent
cover, restricting that cover to the continuous copy of $X$ with
all other coordinates fixed would contradict Lemma~\ref{color}.
\end{proof}

\begin{corollary}
The witness $X$ is locally compact, zero-dimensional, separable,
countably compact and noncompact. Its square $X^2$ is not normal.
\end{corollary}
\begin{proof}
It is an open subspace of the compact Hausdorff zero-dimensional
space $(\beta\N)^2$, giving the first two assertions. The countable
set $\{(a,b)\in\N^2:a\ne b\}$ is dense: around $(p,q)$ with $p\ne q$
choose disjoint basic clopen neighborhoods, and use density of
principal points in each. Countable compactness follows by projection
from $X^\N$. For any free ultrafilter $p$, every neighborhood of
$(p,p)$ contains unequal principal pairs. Thus $X$ is not closed
in $(\beta\N)^2$, and so is not compact.

Finally, $X^2$ is countably compact, hence pseudocompact (every
continuous real-valued function has bounded image). If $X^2$ were
normal, Brandsma--van Hartskamp--van Mill's Theorem~2.1
\cite[p.~28]{BHM} would give a finite functionally open coloring
of every continuous fixed-point-free map on the Tychonoff space $X$.
Applied to $\tau$, this would be a finite open independent cover,
contradicting Lemma~\ref{color}.
\end{proof}

\begin{corollary}\label{general}
For every infinite compact Hausdorff extremally disconnected space $K$,
$D(K)$ has the bounded-clique property and its swap relation has no
countable open independent cover.
\end{corollary}
\begin{proof}
An infinite compact space with a clopen base has countably many disjoint
nonempty clopen sets $V_n$: repeatedly split an infinite clopen remainder
into two nonempty clopen sets, retain an infinite part, and remove the other.
Choose $x_n\in V_n$. For every $A\subseteq\N$, the clopen set
$\cl\bigcup_{n\in A}V_n$ contains $x_n$ exactly when $n\in A$.
The continuous extension $\beta\N\to K$ of $n\mapsto x_n$ is therefore
injective: if two ultrafilters disagree on $A$, their images disagree on
this clopen set. Compactness and Hausdorffness make it an embedding.
Consequently $D(\beta\N)$ embeds in $D(K)$ compatibly with swapping;
a finite open independent cover would restrict to one of $D(\beta\N)$.
Lemma~\ref{power} gives countable compactness and hence also excludes
countable covers. Lemma~\ref{criterion} supplies bounded cliques.
\end{proof}

\begin{corollary}\label{largepower}
Let $I=\beta\N\setminus\N$. Then $X^\N$ is countably compact but
$X^I$ is not. In particular, $X$ is not $u$-compact for any free
ultrafilter $u$ on $\N$.
\end{corollary}
\begin{proof}
Define $z_n\in X^I$ by $z_n(v)=(n,v)$. If $w$ were a cluster point
of $(z_n)$, the second component of $w(v)$ would be $v$ for every $v$.
All first components would equal one point $p\in\beta\N$, since the
equalities between any two first-coordinate projections define closed sets
containing every $z_n$. The point $p$ cannot be a principal natural number
$k$: its singleton clopen neighborhood would meet the sequence only at
$n=k$, contrary to the cluster-point condition. Thus $p\in I$, but then
$w(p)=(p,p)\notin X$, a contradiction.

Here $u$-compact means that every sequence has a $u$-limit. If $X$ were
$u$-compact, the sequence $(z_n)$ would have a coordinatewise $u$-limit
in $X^I$, which is a cluster point because $u$ is free. This proves the
last assertion directly.
\end{proof}

This answers the general-space existence question in
Feldman--Wilce's Remark~4 (p.~8), but that existence assertion already
follows from Ortiz-Castillo, Rodrigues and Tomita~\cite{ORT},
Example~2.3 and Theorem~3.2 (preprint pp.~6--8).
Indeed, their Example~2.3 gives $\N\subseteq Y\subseteq\beta\N$ with
$Y^\kappa$ countably compact for every $\kappa<\mathfrak t$, while
the Vietoris space $\operatorname{CL}(Y)$ of nonempty closed subsets is
not pseudocompact. The tower number $\mathfrak t$ is the least size of a
well-ordered decreasing-modulo-finite family of infinite subsets of $\N$ with no
infinite set almost contained in all its members; $\aleph_0<\mathfrak t$.
If $Y^{\mathfrak c}$ were countably compact, its dense subset
$\N^{\mathfrak c}$ would be relatively countably compact there.
Their Theorem~3.2 would make $\operatorname{CL}(Y)$ pseudocompact,
a contradiction. Here $\mathfrak c=2^{\aleph_0}$, relative countable
compactness means every sequence in the subset has a cluster point in
the ambient space, and pseudocompactness means that every continuous
real-valued function is bounded. Thus their results already imply
$Y^\N$ countably compact and $Y^{\mathfrak c}$ not countably compact.
Corollary~\ref{largepower} is an alternative explicit witness, not a first
existence answer or a topological-group construction.

\begin{corollary}\label{sphere}
The unit sphere $S^2\subseteq\mathbb R^3$, with literal orthogonality,
has $\chi_\op(S^2,\perp)\ge5$.
\end{corollary}
\begin{proof}
Suppose four open orthogonality-free sets $U_1,\ldots,U_4$ cover $S^2$.
A smaller nonempty cover can be enlarged to four by repetitions.
Choose four pairwise disjoint nonempty open balls $B_i\subseteq U_i$.
Assign each point of $B_i$ to class $i$, and assign every other point
to any index whose $U_i$ contains it. This gives an orthogonal partition
$S^2=S_1\sqcup\cdots\sqcup S_4$ with $B_i\subseteq S_i\subseteq U_i$.
By Holmsen--Lee~\cite[Theorem~1, preprint p.~2]{HL}, this partition is
octahedral: in a common orthonormal coordinate frame, each class is
contained in a pair of opposite closed octants. All four distinct pairs
must occur, since the partition covers the sphere. Each $S_i$, and hence
each $U_i$, therefore contains the corresponding pair of open octants.

No $U_i$ can contain a point $q$ of another open octant. After sign changes
and possibly replacing $q$ by $-q$, its coordinates have both signs
relative to the positive octant assigned to $i$. There is a vector $y$
with all coordinates positive and $q\cdot y=0$: choose positive weights
on the positive and negative coordinates to balance the two sums.
Normalize $y$ to lie on $S^2$. It belongs to $U_i$, contradicting
orthogonality-freeness. If $U_i$ contained an octant boundary point,
openness would force it to contain a point in another open octant.
Thus the $U_i$ miss all octant boundaries and cannot cover $S^2$.
\end{proof}

Corollary~\ref{sphere} proves the bound suggested in
Feldman--Wilce's Remark~6. Together with their Section~5 construction
it gives $5\le\chi_\op(S^2,\perp)\le24$; it does not determine the
least number asked for in Question~3.

The following table records source coverage; ``not covered'' describes
the scope of these arguments, not the current status of a problem.
\begin{longtable}{>{\raggedright\arraybackslash}p{.34\linewidth}>{\raggedright\arraybackslash}p{.56\linewidth}}
\hline
Source item & Coverage\\
\hline
Question 2, Section 5, p.~11 & Answered: yes, in ZFC (Theorem~\ref{main});
conditional predecessor:~\cite[Example~2.3]{BHM}.\\
Theorem 1, Section 3, pp.~7--8 & Reproved; prior proof:
Feldman--Wilce (2026), Lemma~\ref{criterion}.\\
Question 1, Section 1, p.~5; selector in Section 4, p.~10 &
Not covered. It asks for failure of the bounded-clique property,
whereas the present witness has that property.\\
Remark 2, p.~5 (necessity of countable compactness for finite cliques;
unbounded cliques without a $T_1$ restriction) & Not covered: the present
witness has bounded cliques. The second question allows $T_1$ witnesses.\\
Remark 4, p.~8 (countable versus higher powers) & Answered: yes;
prior existence proof: Ortiz-Castillo--Rodrigues--Tomita (2017/2018),
Example~2.3 and Theorem~3.2. Alternative witness:
Corollary~\ref{largepower}.\\
Question 3, Section 5, p.~11; Remark 6, p.~12 &
Suggested bound proved (Corollary~\ref{sphere}), using
Holmsen--Lee (2015/2016). Exact value and general-dimensional estimates:
not covered.\\
Countable powers of the present witness & Not posed in the source;
proved in Corollary~\ref{products}.\\
Uncountable swap open chromatic number; infinite extremally disconnected
compact spaces & Not posed; proved in Theorem~\ref{main} and
Corollary~\ref{general}.\\
\cite[Theorem 2.1, p.~28]{BHM} & Prior proof cited; applied above
to deduce that $X^2$ is not normal.\\
\cite[Lemmas 2.1--2.2, pp.~211--212]{PPS} & Prior proof cited:
pair separation and deleted-square countable compactness;
the countable-power consequence is deduced in Section~2.\\
\cite[Question 1.1, p.~27; Question 5.3, p.~39]{BHM} &
Not covered: the present space is not equipped with the required
topological group structure.\\
\hline
\end{longtable}

\section{Verification scope}
The written proofs above cover every assertion of this note.
The following declarations are in the pinned
\href{https://github.com/the-omega-institute/trureturing/blob/eac52692d71b89f89bd7c0d8c734ddc9c5d9e61d/D5/S3/Topology/StoneCech/DeletedSquareChromatic.lean}{\nolinkurl{DeletedSquareChromatic.lean}},
in namespace
\nolinkurl{D5.S3.Topology.StoneCech.DeletedSquareChromatic}:
\begin{itemize}
\item \nolinkurl{countablyCompactSpace_deletedSquare_power}: Lemma~\ref{power}
for compact Hausdorff extremally disconnected spaces;
\item \nolinkurl{boundedCliqueProperty_of_countablyCompact_power}:
the universal graph conclusion of Lemma~\ref{criterion}, including
finiteness of every clique;
\item \nolinkurl{no_finite_open_independent_cover}:
Lemma~\ref{color}, allowing overlapping cover members;
\item \texttt{result}: Hausdorffness of $D(\beta\N)$, its bounded-clique
property and failure of the finite open chromatic property.
\end{itemize}
The clopen agreement statement used for Lemma~\ref{pair} appears as
\nolinkurl{eventually_eq_of_eventually_mem_isClopen_iff} in the pinned
\href{https://github.com/the-omega-institute/trureturing/blob/eac52692d71b89f89bd7c0d8c734ddc9c5d9e61d/D5/S3/Topology/ExtremallyDisconnected/PointPairSeparation.lean}{\nolinkurl{PointPairSeparation.lean}}.
The local-compactness, separability, nonnormal-square, countable-product,
general-space and higher-power corollaries, the uncountable chromatic lower
bound and continuum upper bound, the comparison with~\cite{BHM}, and the
exact swap clique number are written proofs here. The compact
$F$-space countable-power statement is the written deduction from
\cite[Lemma~2.1]{PPS} in Section~2, beyond the extremally disconnected
scope of the displayed Lean declaration. Corollary~\ref{sphere}
is a written consequence of Holmsen--Lee's theorem; the comparison for
Remark~4 uses the cited hyperspace results. The accompanying
\texttt{verify.py} checks finite swap-graph and directional membership-pattern
identities, with the exact ranges in \texttt{README.md} and
\texttt{verification.json}. Those finite checks concern finite graph
identities; the infinite topological statements are proved above.

\begin{thebibliography}{9}
\bibitem{FW}
David V. Feldman and Alexander Wilce,
\emph{A Note on Compactness and Clique Size},
arXiv:2608.13320v1 (2026),
\url{https://arxiv.org/abs/2608.13320v1}.
\bibitem{BHM}
Henno Brandsma, Michael van Hartskamp and Jan van Mill,
\emph{On colorings of topological groups},
Topology Proceedings \textbf{22} (1997), 25--41,
\url{https://staff.fnwi.uva.nl/j.vanmill/papers/papers1997/colorings.pdf}.
\bibitem{PPS}
I. Protasov, J. Pym and D. Strauss,
\emph{A lemma on extending functions into $F$-spaces and homomorphisms
between Stone--\v{C}ech remainders},
Topology and its Applications \textbf{105} (2000), 209--229,
\url{https://doi.org/10.1016/S0166-8641(99)00061-9}.
\bibitem{ORT}
Y. F. Ortiz-Castillo, V. O. Rodrigues and A. H. Tomita,
\emph{Small Cardinals and the Pseudocompactness of Hyperspaces of
Subspaces of $\beta\omega$},
Topology and its Applications \textbf{246} (2018), 9--21;
arXiv:1710.06087v2 (6 November 2017),
\url{https://arxiv.org/abs/1710.06087v2}.
\bibitem{HL}
Andreas F. Holmsen and Seunghun Lee,
\emph{Orthogonal colorings of the sphere},
Mathematika \textbf{62} (2016), 492--501;
arXiv:1505.02514v3 (26 August 2015),
\url{https://arxiv.org/abs/1505.02514v3}.
\end{thebibliography}
\end{document}
